
\documentclass{article}
\usepackage{colortbl}
\usepackage{makecell}
\usepackage{multirow}
\usepackage{supertabular}

\begin{document}

\newcounter{utterance}

\centering \large Interaction Transcript for game `cladder', experiment `full\_v1.5\_default', episode 7517 with qwen3.5{-}4b.
\vspace{24pt}

{ \footnotesize  \setcounter{utterance}{1}
\setlength{\tabcolsep}{0pt}
\begin{supertabular}{c@{$\;$}|p{.15\linewidth}@{}p{.15\linewidth}p{.15\linewidth}p{.15\linewidth}p{.15\linewidth}p{.15\linewidth}}
    \# & \multicolumn{2}{c}{Player} && \multicolumn{2}{c}{Game Master} \\
    \hline

    \theutterance \stepcounter{utterance}  
    & & & \multicolumn{4}{p{0.6\linewidth}}{
        \cellcolor[rgb]{0.9,0.9,0.9}{
            \makecell[{{p{\linewidth}}}]{
                \texttt{\tiny{[P1$\langle$GM]}}
                \texttt{You are an expert in causal inference. The following question is not a typical commonsense query, but rather a meticulously designed question created by a professor specializing in causal inference, intended to assess the students' mastery of the course content.} \\
\texttt{Imagine a self{-}contained, hypothetical world with only the following conditions, and without any unmentioned factors or causal relationships: Wibl has a direct effect on vubr and zuph. Zuph has a direct effect on uvzi. Vubr has a direct effect on uvzi. Method 1: We look at how zuph correlates with uvzi case by case according to vubr. Method 2: We look directly at how zuph correlates with uvzi in general. To understand how zuph affects uvzi, is it more correct to use the Method 1 than Method 2?} \\
\texttt{Start your answer with 'Yes' or 'No', followed by additional reasoning or evidence to support your explanation.} \\
            }
        }
    }
     \\ \\

    \theutterance \stepcounter{utterance}  
    & \multicolumn{4}{p{0.6\linewidth}}{
        \cellcolor[rgb]{0.561,0.737,0.561}{
            \makecell[{{p{\linewidth}}}]{
                \texttt{\tiny{[P1$\rangle$GM]}}
                \texttt{Yes} \\
\\ 
\texttt{**Reasoning:**} \\
\\ 
\texttt{To determine the correct method for estimating the causal effect of \$zuph\$ on \$uvzi\$, we must analyze the causal structure and the presence of confounding variables.} \\
\\ 
\texttt{1.  **Causal Structure Analysis**:} \\
\texttt{    *   \$wibl \textbackslash{}rightarrow zuph\$} \\
\texttt{    *   \$wibl \textbackslash{}rightarrow vubr\$} \\
\texttt{    *   \$vubr \textbackslash{}rightarrow uvzi\$} \\
\texttt{    *   \$zuph \textbackslash{}rightarrow uvzi\$} \\
\\ 
\texttt{    In this diagram, \$wibl\$ is a **common cause** (confounder) of both \$zuph\$ and \$vubr\$. Furthermore, \$vubr\$ is a **collider** relative to the path from \$wibl\$ to \$uvzi\$ if we consider the path \$wibl \textbackslash{}rightarrow vubr \textbackslash{}rightarrow uvzi\$, but more importantly, \$vubr\$ lies on a backdoor path from \$zuph\$ to \$uvzi\$ if we do not condition on it correctly. Specifically, the path \$zuph \textbackslash{}leftarrow wibl \textbackslash{}rightarrow vubr \textbackslash{}rightarrow uvzi\$ is a backdoor path that introduces spurious correlation between \$zuph\$ and \$uvzi\$ unless blocked.} \\
\\ 
\texttt{2.  **Evaluating Method 2 (General Correlation)**:} \\
\texttt{    Method 2 looks at the raw correlation between \$zuph\$ and \$uvzi\$. This correlation is biased because it includes the effect of the common cause \$wibl\$. Since \$wibl\$ affects both \$zuph\$ and \$vubr\$, and \$vubr\$ affects \$uvzi\$, the variation in \$uvzi\$ driven by \$wibl\$ (via \$vubr\$) will be correlated with \$zuph\$ (also driven by \$wibl\$). Therefore, Method 2 confounds the direct causal effect of \$zuph\$ on \$uvzi\$ with the indirect effect flowing through \$vubr\$ (which is also driven by \$wibl\$). This is a classic case of **confounding bias**.} \\
\\ 
\texttt{3.  **Evaluating Method 1 (Stratification/Conditioning)**:} \\
\texttt{    Method 1 looks at the correlation between \$zuph\$ and \$uvzi\$ "case by case according to \$vubr\$." This implies conditioning on \$vubr\$ (stratifying the data by the levels of \$vubr\$).} \\
\texttt{    *   By conditioning on \$vubr\$, we block the backdoor path \$zuph \textbackslash{}leftarrow wibl \textbackslash{}rightarrow vubr \textbackslash{}rightarrow uvzi\$.} \\
\texttt{    *   Once \$vubr\$ is fixed (conditioned upon), the influence of \$wibl\$ on \$uvzi\$ is severed for that specific stratum.} \\
\texttt{    *   Consequently, the remaining correlation between \$zuph\$ and \$uvzi\$ within each stratum of \$vubr\$ represents the direct causal effect of \$zuph\$ on \$uvzi\$, free from the confounding influence of \$wibl\$.} \\
\\ 
\texttt{4.  **Conclusion**:} \\
\texttt{    Since Method 2 suffers from confounding bias due to the common cause \$wibl\$, it yields an incorrect estimate of the causal effect. Method 1 correctly blocks the backdoor path by conditioning on the mediator/confounder \$vubr\$, thereby isolating the true causal effect. Therefore, Method 1 is the more correct approach for understanding how \$zuph\$ affects \$uvzi\$.} \\
            }
        }
    }
    & & \\ \\

    \theutterance \stepcounter{utterance}  
    & & & \multicolumn{2}{p{0.3\linewidth}}{
        \cellcolor[rgb]{0.9,0.9,0.9}{
            \makecell[{{p{\linewidth}}}]{
                \texttt{\tiny{[GM$|$GM]}}
                \texttt{Yes} \\
            }
        }
    }
    & & \\ \\

    \theutterance \stepcounter{utterance}  
    & & & \multicolumn{2}{p{0.3\linewidth}}{
        \cellcolor[rgb]{0.9,0.9,0.9}{
            \makecell[{{p{\linewidth}}}]{
                \texttt{\tiny{[GM$|$GM]}}
                \texttt{yes} \\
            }
        }
    }
    & & \\ \\

    \theutterance \stepcounter{utterance}  
    & & & \multicolumn{2}{p{0.3\linewidth}}{
        \cellcolor[rgb]{0.9,0.9,0.9}{
            \makecell[{{p{\linewidth}}}]{
                \texttt{\tiny{[GM$|$GM]}}
                \texttt{game\_result = WIN} \\
            }
        }
    }
    & & \\ \\

\end{supertabular}
}

\end{document}
